\section{Test 3: Proporcionalitat i aplicacions}

\iniciTest{tC}{c,b,d,a,c,b,d,c,b,d,b,d}

\begin{enumerate}
\pregunta Amb \(3\) litres de pintura es poden pintar \(24\) m\(^2\) de paret. Amb \(5\) litres, mantenint el mateix rendiment, es podran pintar:
\begin{enumerate}
\opcio{a}{\(30\) m\(^2\).}
\opcio{b}{\(36\) m\(^2\).}
\opcio{c}{\(40\) m\(^2\).}
\opcio{d}{\(45\) m\(^2\).}
\end{enumerate}
\feedback

\pregunta Quatre treballadors fan una feina en \(15\) dies. Sis treballadors, treballant al mateix ritme, trigaran:
\begin{enumerate}
\opcio{a}{\(8\) dies.}
\opcio{b}{\(10\) dies.}
\opcio{c}{\(12\) dies.}
\opcio{d}{\(22{,}5\) dies.}
\end{enumerate}
\feedback

\pregunta Tres persones han dedicat \(2\), \(3\) i \(4\) hores, respectivament, a una mateixa activitat. Si es reparteixen \(900\) euros segons el temps dedicat, les parts són:
\begin{enumerate}
\opcio{a}{\(300\), \(300\), \(300\).}
\opcio{b}{\(150\), \(300\), \(450\).}
\opcio{c}{\(400\), \(300\), \(200\).}
\opcio{d}{\(200\), \(300\), \(400\).}
\end{enumerate}
\feedback

\pregunta Tres màquines fan una mateixa tasca en \(2\), \(3\) i \(6\) hores. Es vol repartir una prima de \(780\) euros premiant més la màquina més ràpida. Les parts han de ser:
\begin{enumerate}
\opcio{a}{\(390\), \(260\), \(130\).}
\opcio{b}{\(130\), \(260\), \(390\).}
\opcio{c}{\(260\), \(260\), \(260\).}
\opcio{d}{\(300\), \(260\), \(220\).}
\end{enumerate}
\feedback

\pregunta En una botiga, \(3\) kg d'arròs costen \(7{,}20\) euros. En una altra botiga, \(5\) kg costen \(11{,}50\) euros. Quina opció és més barata per quilogram?
\begin{enumerate}
\opcio{a}{La primera botiga, perquè \(7{,}20<11{,}50\).}
\opcio{b}{Les dues tenen el mateix preu per quilogram.}
\opcio{c}{La segona botiga.}
\opcio{d}{No es pot saber sense comprar la mateixa quantitat.}
\end{enumerate}
\feedback

\pregunta Una recepta necessita \(250\) g de farina per a \(4\) racions. Per a \(10\) racions calen:
\begin{enumerate}
\opcio{a}{\(500\) g.}
\opcio{b}{\(625\) g.}
\opcio{c}{\(750\) g.}
\opcio{d}{\(1000\) g.}
\end{enumerate}
\feedback

\pregunta Una beguda es prepara amb \(2\) parts de concentrat i \(5\) parts d'aigua. Per preparar \(3{,}5\) litres de beguda, calen:
\begin{enumerate}
\opcio{a}{\(0{,}5\) L de concentrat i \(3\) L d'aigua.}
\opcio{b}{\(1{,}5\) L de concentrat i \(2\) L d'aigua.}
\opcio{c}{\(2\) L de concentrat i \(1{,}5\) L d'aigua.}
\opcio{d}{\(1\) L de concentrat i \(2{,}5\) L d'aigua.}
\end{enumerate}
\feedback

\pregunta En un plànol a escala \(1:200\), una habitació fa \(4\) cm. La longitud real és:
\begin{enumerate}
\opcio{a}{\(4\) m.}
\opcio{b}{\(6\) m.}
\opcio{c}{\(8\) m.}
\opcio{d}{\(20\) m.}
\end{enumerate}
\feedback

\pregunta Quin dels punts següents pertany a la funció \(y=2x\)?
\begin{enumerate}
\opcio{a}{\((2,5)\).}
\opcio{b}{\(\left(\frac{3}{2},3\right)\).}
\opcio{c}{\((3,5)\).}
\opcio{d}{\(\left(\frac{1}{2},2\right)\).}
\end{enumerate}
\feedback

\pregunta Quin dels punts següents pertany a la funció \(y=\frac{6}{x}\)?
\begin{enumerate}
\opcio{a}{\((2,2)\).}
\opcio{b}{\((3,3)\).}
\opcio{c}{\((6,6)\).}
\opcio{d}{\(\left(4,\frac{3}{2}\right)\).}
\end{enumerate}
\feedback

\pregunta Considerem la taula
\[
\begin{array}{c|ccc}
x & 2 & 3 & 6 \\
\hline
y & 12 & 8 & 4
\end{array}
\]
La relació entre \(x\) i \(y\) és:
\begin{enumerate}
\opcio{a}{De proporcionalitat directa.}
\opcio{b}{De proporcionalitat inversa.}
\opcio{c}{Les dues coses alhora.}
\opcio{d}{Cap relació de proporcionalitat.}
\end{enumerate}
\feedback

\pregunta Considerem la taula
\[
\begin{array}{c|ccc}
x & 1 & 2 & 3 \\
\hline
y & 3 & 5 & 7
\end{array}
\]
La relació entre \(x\) i \(y\) és:
\begin{enumerate}
\opcio{a}{De proporcionalitat directa.}
\opcio{b}{De proporcionalitat inversa.}
\opcio{c}{Una relació de proporcionalitat directa amb constant \(3\).}
\opcio{d}{Cap relació de proporcionalitat.}
\end{enumerate}
\feedback
\end{enumerate}
